\capitulo{7}{Conclusiones}
La ejecución del presente Trabajo de Fin de Grado ha culminado con el desarrollo y despliegue de un \textit{software} innovador para la gestión, monitorización y cálculo de indicadores clínicos en la unidad de Reanimación (\textit{URCCPQ}), cumpliendo y superando los objetivos iniciales. 

El inicio de este proyecto supuso un reto de ingeniería de requisitos considerable. Al estar dirigido a facultativos médicos, cuyo perfil carece de formación técnica específica en informática, la problemática se planteó inicialmente de forma sumamente abierta, sin una arquitectura predefinida ni limitaciones tecnológicas claras. Esta abstracción inicial obligó a asumir el liderazgo del diseño desde su concepción, analizando de forma autónoma qué pila tecnológica (\textit{stack}) resultaría idónea para ofrecer una solución robusta y, sobre todo, accesible.

La verdadera naturaleza de este proyecto trasciende el mero desarrollo de \textit{software}, justificando plenamente su pertinencia y valor dentro del ámbito de la\textbf{ Ingeniería de la Salud} frente a la Ingeniería Informática tradicional. Mientras que un enfoque puramente informático prioriza la optimización algorítmica y la arquitectura de los datos, el ingeniero de la salud actúa como el puente indispensable entre la tecnología y la práctica clínica. 

El diseño de esta plataforma ha exigido ir mucho más allá del código: ha requerido una comprensión profunda del lenguaje médico, de la criticidad de variables de soporte vital (como la escala \textit{APACHE}, el \textit{shock} séptico o los días de \textit{VMI}) y de los exigentes protocolos de una \textbf{Unidad de Reanimación}. Esta capacidad para comprender el entorno clínico, empatizar con el flujo de trabajo crítico del facultativo bajo presión y traducir necesidades médicas complejas en soluciones tecnológicas usables, seguras y rigurosas con la privacidad del paciente es, en esencia, el valor diferencial que aporta la \textbf{Ingeniería de la Salud} a la sociedad.

A pesar de la incertidumbre inherente a las fases iniciales, la libertad metodológica para proponer y desarrollar una arquitectura de \textit{software} propia ha consolidado una profunda confianza en la capacidad de resolución de problemas complejos. Resulta altamente satisfactorio haber logrado transformar una necesidad clínica abstracta en una herramienta informática funcional y de grado de producción. El sistema final optimiza y simplifica drásticamente el procesamiento de datos en la unidad, garantizando un acabado profesional y sentando las bases para futuras implementaciones hacia un ecosistema integral de monitorización hospitalaria.

\section{Aspectos relevantes y desafíos técnicos}

Durante el ciclo de vida del desarrollo se presentaron diversos obstáculos técnicos que requirieron la reevaluación y adaptación de las estrategias iniciales:

El primer gran desafío radicó en la fase de investigación y selección del entorno de desarrollo. La necesidad de aunar un potente motor estadístico con una interfaz \textit{web} moderna e interactiva generó dificultades a la hora de determinar qué lenguaje y \textit{framework} emplear. Tras evaluar distintas alternativas, se optó por unificar el desarrollo bajo \textit{Python} empleando \textit{Reflex}. Esta decisión implicó enfrentarse a la curva de aprendizaje de tecnologías de despliegue novedosas, pero resultó vital para superar las limitaciones estructurales y visuales que presentaban otras opciones más tradicionales de la ciencia de datos.

El segundo aspecto crítico, y uno de los mayores escollos del proyecto, tuvo lugar durante la fase de despliegue local. Con el objetivo de facilitar el uso a los doctores, el planteamiento original buscaba compilar y empaquetar toda la aplicación en un único archivo ejecutable tradicional (\textit{.exe}). Sin embargo, la complejidad arquitectónica de un \textit{framework full-stack}\footnote{Requiere levantar servidores \textit{FastAPI} y compilar entornos \textit{JavaScript} simultáneamente.} hizo descartar esta alternativa por incompatibilidades arquitectónicas críticas fundamentadas en los siguientes factores:

\begin{itemize}
	\item \textbf{Arquitectura distribuida y WebSockets:} \textit{Reflex} opera bajo un modelo \textit{full-stack} que despliega dos servidores simultáneos (\textit{Frontend} y \textit{Backend}). La sincronización del estado de la aplicación y la transferencia de datos entre ambos dependen de una conexión bidireccional continua a través de \textit{WebSockets}.
	
	\item \textbf{Ejecución en directorios temporales:} El empaquetado en un único archivo obliga a que el código se descomprima dinámicamente en una carpeta temporal del sistema anfitrión (como el directorio \textit{\_MEI} en la ruta \textit{AppData} de \textit{Windows}) cada vez que se ejecuta.
	
	\item \textbf{Restricciones de seguridad del sistema operativo:} Los \textit{firewalls} y sistemas de protección nativos (como \textit{Windows Defender}) aplican políticas estrictas sobre los procesos iniciados desde rutas temporales. Al intentar levantar servidores locales y abrir puertos de red desde estas ubicaciones, el sistema operativo interpreta la acción como una posible amenaza, bloqueando silenciosamente el tráfico de los \textit{WebSockets}.
\end{itemize}

Como resultado de este bloqueo, la interfaz gráfica lograba renderizarse en el navegador, pero la comunicación con la lógica del servidor (despliegue de eventos, subida de archivos) quedaba completamente anulada. Como solución de ingeniería, se rediseñó la estrategia de despliegue mediante la orquestación de contenedores \textit{Docker} directamente en la \textit{intranet} del hospital, garantizando el correcto flujo de red y el estricto respeto a las políticas de seguridad institucionales. Esta alternativa superó los problemas de compatibilidad del entorno virtual, proporcionando a los facultativos exactamente la misma experiencia de acceso directo e inmediato (\textit{plug and play}) que se buscaba inicialmente, concluyendo así el trabajo con el cumplimiento íntegro de las expectativas.

Finalmente, un tercer desafío relevante fue la evolución del sistema de gestión de la información. El planteamiento original contemplaba una persistencia basada exclusivamente en la carga manual de archivos planos de tipo \textit{CSV}. No obstante, con el fin de profesionalizar la herramienta, garantizar la integridad de los registros y permitir un control de usuarios riguroso, se optó por la implementación obligatoria de una base de datos relacional. Para ello, se seleccionó \textit{PostgreSQL} en su versión 15 o superior, decisión motivada por la necesidad de asegurar la compatibilidad nativa con la infraestructura tecnológica ya operativa en el \textit{HUBU}. Este cambio arquitectónico permitió abandonar el manejo de archivos locales aislados en favor de una solución integrada, capaz de soportar el histórico clínico de la unidad bajo los estándares de seguridad y rendimiento exigidos por el hospital.



